\documentclass[12pt,a4paper]{article}

% ── Packages ────────────────────────────────────────────────────────────────────
\usepackage[utf8]{inputenc}
\usepackage[T1]{fontenc}
\usepackage{lmodern}
\usepackage[margin=2.5cm]{geometry}
\usepackage{amsmath,amssymb}
\usepackage{graphicx}
\usepackage{booktabs}
\usepackage{multirow}
\usepackage{array}
\usepackage{hyperref}
\usepackage{xcolor}
\usepackage{natbib}
\usepackage{caption}
\usepackage{subcaption}
\usepackage{microtype}
\usepackage{setspace}
\usepackage{parskip}
\usepackage{fancyhdr}
\usepackage{titlesec}
\usepackage{enumitem}

% ── Hyperref config ─────────────────────────────────────────────────────────────
\hypersetup{
  colorlinks=true,
  linkcolor=blue!60!black,
  citecolor=green!50!black,
  urlcolor=blue!70!black,
  pdftitle={Do Celestial Bodies Influence Tropical Cyclone Behaviour?},
  pdfauthor={Anonymous}
}

% ── Header / footer ─────────────────────────────────────────────────────────────
\pagestyle{fancy}
\fancyhf{}
\rhead{\small Tropical Cyclones \& Vedic Astronomy}
\lhead{\small EDA Pipeline Analysis (2016--2026)}
\cfoot{\thepage}

% ── Title ───────────────────────────────────────────────────────────────────────
\title{
  {\Large \textbf{Do Celestial Bodies Influence Tropical Cyclone Behaviour?}}\\[0.4em]
  {\large A Rigorous Exploratory Analysis of Vedic Astronomical and Panchang Features\\
  Against a Decade of Global Tropical Cyclone Track Data}
}
\author{
  Anonymous Submission\\[0.3em]
  \small \textit{Empirical investigation --- not intended as endorsement of astrological claims}
}
\date{\today}

\begin{document}

\maketitle
\thispagestyle{fancy}

\begin{abstract}
\noindent
We present a large-scale exploratory data analysis examining whether features derived
from Vedic Panchang (Hindu lunar calendar) and classical astronomical ephemeris
calculations carry statistically significant associations with tropical cyclone
intensity and behaviour.
Our dataset spans 944 tropical cyclones observed between July 2016 and January 2026,
producing 47\,533 track observations enriched with 282 derived features covering
standard meteorological quantities, Skyfield-computed planetary positions, and
Vedic categorical variables (Tithi, Nakshatra, Yoga, Karana, Paksha, Samvatsara, etc.).

We apply a seven-stage analysis pipeline: (1) data transformation and feature
engineering, (2) basic exploratory analysis with multiple-comparison correction,
(3) advanced feature importance ranking with storm-grouped cross-validation,
(4) SHAP/LIME explainability, (5) comprehensive lifecycle and circular-statistics
analysis, (6) formal hypothesis testing, and (7) block-wise predictive modelling.

\textbf{Key findings:}
(i) Adding Vedic features to a meteorological baseline \emph{degrades} out-of-sample
    RMSE by 12.6\% (baseline 4.04~kt vs.\ Vedic 4.55~kt), indicating that Vedic
    features add noise rather than predictive signal.
(ii) A marginal association between lunar phase and rapid intensification
    ($\chi^2 = 19.66$, $p = 0.0064$, 7 df) survives Bonferroni correction
    ($\alpha_{\text{adj}} = 0.0071$), but the effect size is small and the physical
    mechanism is unclear.
(iii) Samvatsara-year ANOVA reaches $p < 10^{-130}$, but this reflects known
    interannual climate variability (ENSO/PDO), not Hindu calendar cycles.
(iv) Storm-grouped cross-validation reveals that all tested tree-ensemble models
    achieve negative R\textsuperscript{2} ($-0.09$ to $-0.11$), confirming that the
    astronomical feature block provides no out-of-storm generalisation.
(v) Geospatial and track-geometry features (\textit{latitude}, \textit{longitude},
    \textit{spline\_bearing\_deg}) remain the dominant predictors throughout, as
    expected from tropical cyclone climatology.

We discuss methodological pitfalls --- cyclical-variable treatment, within-storm
autocorrelation leakage, and the Samvatsara confound --- and provide corrected
analyses for each. Our conclusion is that the current evidence does not support
a predictive or associative role for Vedic astronomical constructs in tropical
cyclone intensity.
\end{abstract}

\tableofcontents
\newpage

% ═══════════════════════════════════════════════════════════════════════════════
\section{Introduction}
% ═══════════════════════════════════════════════════════════════════════════════

Tropical cyclones are among the most destructive natural phenomena on Earth,
responsible for billions of dollars in damage and thousands of casualties annually.
Accurate intensity prediction remains an open research challenge: despite advances
in numerical weather prediction and statistical modelling, the 24--48~hour intensity
forecast skill of state-of-the-art models plateaus because of complex multi-scale
interactions involving sea surface temperature (SST), vertical wind shear, moisture
fields, and internal dynamics~\citep{demaria2014advances,zhang2018tropical}.

An independent thread of inquiry asks whether \emph{astronomical} or
\emph{astrological} quantities --- planetary positions, lunar phase, Vedic calendar
cycles --- are associated with cyclone behaviour. Such a hypothesis is physically
speculative: gravity from the Moon and Sun produce tidal stresses on the ocean and
atmosphere, but the magnitudes involved are many orders of magnitude below the
energy budget of a mature tropical cyclone.
Nevertheless, the hypothesis is empirically testable, and a rigorous null-result
analysis can clarify the upper-bound size of any such associations, guiding
future data collection.

This work assembles the first large-scale dataset merging tropical cyclone track
archives with Vedic Panchang / Jyotisha-Vedanga astronomical quantities. We subject
this dataset to a carefully controlled analysis pipeline with explicit
multiple-comparison correction, storm-grouped cross-validation, and circular
statistics for angular planetary variables. We report results transparently,
including negative or null findings.

\subsection{Research Questions}

\begin{enumerate}[label=\textbf{RQ\arabic*.}]
  \item Do Vedic Panchang features improve out-of-sample prediction of cyclone
        wind intensity beyond a meteorological baseline?
  \item Is there a statistically robust association between lunar phase and the
        rate of rapid intensification (RI) events?
  \item Do planetary speed signatures (direct vs.\ retrograde motion) correlate
        with cyclone trajectory dynamics?
  \item Does the Vedic Samvatsara year cycle explain variance in storm intensity
        beyond interannual meteorological variability?
  \item Which astronomical or Vedic features show the most consistent importance
        across multiple analysis methods (tree models, SHAP, LIME, mutual information)?
\end{enumerate}

% ═══════════════════════════════════════════════════════════════════════════════
\section{Data}
% ═══════════════════════════════════════════════════════════════════════════════

\subsection{Source Data}

\textbf{Cyclone track data.} We use 944~JSON files from a proprietary storm archive
covering the period 2016-07-15 to 2026-01-24~UTC. Each file describes a single
tropical cyclone with timestamped position, wind speed (kt), minimum central pressure
(hPa), storm type classification, and advisory metadata. After join and deduplication,
the dataset contains 47\,533 track observations representing 735 distinct storms
(a subset of the 944 files that yielded usable track rows after merging).

\textbf{Astronomical ephemeris data.} For each track observation, geocentric and
topocentric positions for nine solar-system bodies
(Sun, Moon, Mercury, Venus, Mars, Jupiter, Saturn, Uranus, Neptune) are computed
using the Skyfield Python library~\citep{skyfield2019}.
Per-body fields include right ascension, declination, ecliptic longitude and latitude,
altitude and azimuth above the local horizon, and sin/cos-transformed variants.

\textbf{Vedic Panchang data.} Vedic calendar attributes are extracted from a
Drig-Ganita (observational) Panchang computation engine stored alongside each
track record. These include:
\begin{itemize}
  \item \textbf{Tithi} (lunar day, 1--30), \textbf{Nakshatra} (lunar mansion, 1--27),
        \textbf{Yoga} (luni-solar yoga, 1--27), \textbf{Karana} (half-tithi, 1--11),
        \textbf{Paksha} (fortnight: Shukla/Krishna), \textbf{Moonsign} (Rashi, 1--12),
        \textbf{Sunsign} (solar rashi), \textbf{Drik Ritu} (season)
  \item \textbf{Samvatsara} (60-year Jupiter cycle), \textbf{Drik Ayana} (solstice half)
  \item Numeric Vedic planetary quantities: full degree (0--360\textdegree),
        padam (nakshatra subdivision), speed (degrees/day), right ascension, declination
        for eight classical Vedic grahas (Surya, Chandra, Mangal, Budha, Guru, Shukra,
        Shani, Rahu/Ketu).
\end{itemize}

\textbf{Derived features.} The transformation pipeline computes:
\begin{itemize}
  \item Temporal: \textit{hour\_utc}, \textit{day\_of\_week}, \textit{day\_of\_year}, \textit{month}
  \item Storm dynamics: \textit{wind\_change} (kt/step), \textit{pressure\_change} (hPa/step),
        \textit{displacement} (degrees), \textit{movement\_speed} (deg/h)
  \item Spline geometry: smoothed forecast track bearing, curvature, length, and endpoint
        (\textit{spline\_bearing\_deg}, \textit{spline\_curvature}, etc.)
  \item Day-length: \textit{dinamana\_minutes} (daylight duration), \textit{ratrimana\_minutes}
\end{itemize}

\subsection{Dataset Summary}

\begin{table}[h]
\centering
\caption{Dataset overview after transformation and join.}
\label{tab:dataset}
\begin{tabular}{lr}
\toprule
\textbf{Property} & \textbf{Value} \\
\midrule
Observations (rows)               & 47\,533 \\
Feature columns                   & 282 \\
Unique storms                     & 735 \\
Date range                        & 2016-07-15 to 2026-01-24 UTC \\
Wind speed: mean / median (kt)    & 49.9 / 40.0 \\
Wind speed: SD / range (kt)       & 27.4 / [0, 170] \\
Pressure: mean (hPa)              & 989.9 (SD = 19.5) \\
Latitude range (\textdegree)      & $[-44.3,\ 64.0]$ \\
Longitude range (\textdegree)     & $[-179.8,\ 256.3]$ \\
Astronomical (Skyfield) features  & 90 \\
Vedic planetary features          & 75 \\
Vedic categorical features        & 10 \\
Meteorological features           & 14 \\
\bottomrule
\end{tabular}
\end{table}

\subsection{Missing Data}

Table~\ref{tab:missing} summarises columns with $> 5\%$ missing values.
The \textit{storm\_active} flag (99.7\% null) and cone forecast features
($> 99\%$ null) are essentially absent and are excluded from all models.
Spline features ($\approx 15.5\%$ null) are retained with listwise deletion
where used.

\begin{table}[h]
\centering
\caption{Columns with more than 5\% missing values.}
\label{tab:missing}
\begin{tabular}{lc}
\toprule
\textbf{Feature Group} & \textbf{Missing (\%)} \\
\midrule
\textit{storm\_active} flag            & 99.7 \\
Cone geometry features (7 columns)     & $\approx 99.0$ \\
Spline geometry features (7 columns)   & $\approx 15.5$ \\
Anomaly label / score                  & 10.9 \\
\bottomrule
\end{tabular}
\end{table}

% ═══════════════════════════════════════════════════════════════════════════════
\section{Analysis Pipeline}
% ═══════════════════════════════════════════════════════════════════════════════

The analysis proceeds through seven modular stages implemented as independent
Python scripts. Each stage was audited for methodological correctness and
corrected where necessary; the corrections are described in Section~\ref{sec:corrections}.

\begin{enumerate}
  \item \textbf{Data transformation} (\texttt{01\_data\_transform.py}): JSON ingestion,
        Skyfield ephemeris computation, Vedic quantity extraction, feature engineering.
  \item \textbf{Basic EDA} (\texttt{02\_basic\_eda.py}): Descriptive statistics,
        Kruskal--Wallis tests with Benjamini--Hochberg FDR correction,
        correlation screening.
  \item \textbf{Advanced EDA} (\texttt{03\_advanced\_eda.py}): PCA, tree-ensemble
        feature importance with storm-\emph{grouped} cross-validation, mutual
        information, clustering, redundancy analysis.
  \item \textbf{SHAP/LIME} (\texttt{04\_shap\_lime.py}): Global and local model
        explanation with consensus feature identification.
  \item \textbf{Comprehensive EDA} (\texttt{05\_comprehensive\_eda.py}): Lifecycle
        analysis, circular-linear correlations for angular features, retrograde
        analysis with significance tests.
  \item \textbf{Hypothesis testing} (\texttt{06\_hypothesis\_testing.py}): Formal
        statistical tests for Vedic predictive lift (H1), lunar-phase RI
        association (H2), planetary speed -- trajectory coupling (H3),
        Samvatsara intensity analysis, Mars/Rahu/Ketu analysis.
  \item \textbf{Predictive modelling} (\texttt{07\_predictive\_model\_stack.py}):
        Block-wise Random Forest regression comparing meteorology-only, panchang-enhanced,
        and full feature stacks using temporal forward-chaining cross-validation.
\end{enumerate}

% ═══════════════════════════════════════════════════════════════════════════════
\section{Methodological Corrections}
\label{sec:corrections}
% ═══════════════════════════════════════════════════════════════════════════════

During the pipeline audit we identified and corrected five methodological issues.
We document them here for full transparency.

\subsection{Circular-Variable Treatment (E2)}

Planetary ecliptic longitude and Vedic full-degree quantities are angular variables
$\theta \in [0\textdegree, 360\textdegree)$.
Applying Pearson $r$ directly treats the interval as linear and incorrectly penalises
the 0\textdegree/360\textdegree discontinuity (e.g., $\theta = 359\textdegree$ is
treated as maximally far from $\theta = 1\textdegree$).
We corrected this by replacing Pearson $r$ with a circular--linear decomposition:
\begin{equation}
  r_{\sin} = \mathrm{Corr}\bigl(\sin(\theta),\, y\bigr), \quad
  r_{\cos} = \mathrm{Corr}\bigl(\cos(\theta),\, y\bigr), \quad
  r_{\text{circ}} = \max\!\bigl(|r_{\sin}|,\, |r_{\cos}|\bigr)
\end{equation}
Results reported in \texttt{vedic\_degree\_correlations.csv} now use $r_{\text{circ}}$.

\subsection{Within-Storm Cross-Validation Leakage (C3)}

The original 5-fold CV split rows randomly, allowing observations from the same storm
to appear in both training and test folds. Because consecutive track points of a storm
are strongly autocorrelated, this inflates $R^2$ estimates by effectively testing
memorisation rather than generalisation. We replaced standard \texttt{KFold} with
\texttt{GroupKFold(n\_splits=5)} grouped by \texttt{storm\_id}, ensuring all
observations from a given storm appear in exactly one fold.

\subsection{Multiple-Comparison Correction (B1)}

The Kruskal--Wallis tests in Stage 2 test 10 Vedic category variables simultaneously.
At $\alpha = 0.05$, approximately 0.5 false discoveries are expected by chance.
We applied Benjamini--Hochberg (BH) FDR correction and added columns
\texttt{p\_adjusted\_bh} and \texttt{significant\_after\_bh\_correction} to
\texttt{kruskal\_wallis\_wind.csv}. All conclusions use BH-adjusted $p$-values.

\subsection{R\textsuperscript{2} Discrepancy Between Stages 3 and 4}

Stage~3 (grouped CV) reports $R^2 \approx -0.10$, while Stage~4 (SHAP on XGBoost)
reports test $R^2 = 0.71$. The discrepancy arises from three sources:
(1)~Stage~4 uses a stricter $> 80\%$ completeness filter retaining only the most
complete features; (2)~Stage~4 uses a random 80/20 split, permitting within-storm
autocorrelation to inflate the estimate; (3)~Stage~3 uses GroupKFold.
The grouped-CV $R^2$ is the appropriate measure of out-of-storm generalisation.

\subsection{Rahu / Ketu Retrograde Independence}

Rahu (North Lunar Node) and Ketu (South Lunar Node) are always exactly antiopodal on
the ecliptic and always move together at identical speed magnitude (retrograde
$\approx -0.053\textdegree$/day). Any statistical test run on both independently will
produce \emph{identical} results by astronomical definition; this is not a
copy-paste error but an inherent constraint. Results for Ketu are labelled accordingly.

% ═══════════════════════════════════════════════════════════════════════════════
\section{Results}
% ═══════════════════════════════════════════════════════════════════════════════

\subsection{Descriptive Statistics and Data Quality}

Wind speed follows a right-skewed distribution (mean = 49.9~kt, median = 40.0~kt,
SD = 27.4~kt, skewness = 1.35, range = [0, 170]~kt). Central pressure is
approximately symmetric (mean = 989.9~hPa, SD = 19.5~hPa, range = [882, 1019]~hPa).
Wind and pressure show a strong negative correlation ($r^2 = 0.886$).

The 208 numeric features can be grouped as: 90 Skyfield astronomical features,
75 Vedic planetary features, and a smaller meteorological/derived set.
PCA on the 90-dimensional Skyfield feature matrix shows that 29 principal components
are needed to explain $\geq 90\%$ of variance, confirming moderate redundancy
but non-trivial intrinsic dimensionality consistent with the planet-body structure
of the feature space (PC1 is dominated by Jupiter--Uranus--Neptune ecliptic longitude
cluster; PC2 captures Sun/inner-planet altitude/azimuth patterns).
A total of 239 feature pairs exhibit $|r| > 0.95$, predominantly among
the sin/cos transforms of the same body position.

\subsection{Kruskal--Wallis Tests on Vedic Categories (BH-Corrected)}
\label{sec:kw}

Table~\ref{tab:kw} reports Kruskal--Wallis H statistics and BH-corrected $p$-values
for ten Vedic categorical variables against wind speed.
Eight of ten categories are significant after correction, with Drik Ritu
(season: $H = 1214$) and Sunsign (solar zodiac: $H = 1123$) showing the largest
effect. \emph{However}, as noted in Section~\ref{sec:confound}, Ritu and Sunsign are
essentially time-of-year proxies that encode the annual tropical cyclone season cycle.
Weekday is not significant (H = 11.6, $p = 0.072$).

\begin{table}[h]
\centering
\caption{Kruskal--Wallis tests (wind speed by Vedic category) with BH correction.
  ``Sig.''\ = significant at BH-adjusted $p < 0.05$.}
\label{tab:kw}
\begin{tabular}{lrrrc}
\toprule
\textbf{Category} & \textbf{$H$} & \textbf{$p$ (raw)} & \textbf{$p$ (BH adj.)} & \textbf{Sig.} \\
\midrule
Drik Ritu (season)   & 1214.0  & $2.7\times10^{-260}$ & $2.7\times10^{-259}$ & Yes \\
Sunsign              & 1123.1  & $5.9\times10^{-234}$ & $2.9\times10^{-233}$ & Yes \\
Yoga                 & 143.3   & $3.4\times10^{-18}$  & $8.6\times10^{-18}$  & Yes \\
Nakshatra            & 116.1   & $2.4\times10^{-13}$  & $4.7\times10^{-13}$  & Yes \\
Drik Ayana           & 102.6   & $4.0\times10^{-24}$  & $1.3\times10^{-23}$  & Yes \\
Moonsign             & 77.9    & $3.7\times10^{-12}$  & $6.2\times10^{-12}$  & Yes \\
Tithi                & 51.3    & $7.5\times10^{-6}$   & $9.4\times10^{-6}$   & Yes \\
Karana               & 42.5    & $6.2\times10^{-6}$   & $8.8\times10^{-6}$   & Yes \\
Paksha (fortnight)   & 8.1     & $4.4\times10^{-3}$   & $4.9\times10^{-3}$   & Yes \\
Weekday              & 11.6    & $7.2\times10^{-2}$   & $7.2\times10^{-2}$   & No  \\
\bottomrule
\end{tabular}
\end{table}

\subsection{Feature Importance under Storm-Grouped Cross-Validation}
\label{sec:importance}

Using \texttt{GroupKFold} (groups = storm\_id), the RF, XGBoost, and LightGBM models
trained on the full 186-feature set all show negative grouped CV $R^2$:

\begin{center}
\begin{tabular}{lcc}
\toprule
\textbf{Model} & \textbf{Grouped CV $R^2$} & \textbf{(SD)} \\
\midrule
Random Forest  & $-0.0901$ & $\pm 0.0397$ \\
XGBoost        & $-0.1089$ & $\pm 0.0500$ \\
LightGBM       & $-0.1010$ & $\pm 0.0262$ \\
\bottomrule
\end{tabular}
\end{center}

Negative $R^2$ means the model performs worse than predicting the global mean wind for
every new (unseen) storm.
Despite this, model-internal feature importances indicate which features the model
\emph{uses} (not which are predictive in a generalisable sense).
The top 10 averaged importances are shown in Table~\ref{tab:importance}.

\begin{table}[h]
\centering
\caption{Top-10 averaged feature importances (RF/XGB/LGB mean importance score).
  Note: models show negative grouped CV $R^2$; rankings are exploratory, not predictive.}
\label{tab:importance}
\begin{tabular}{lc}
\toprule
\textbf{Feature} & \textbf{Mean importance} \\
\midrule
\textit{longitude}                        & 0.0588 \\
\textit{latitude}                         & 0.0420 \\
\textit{spline\_end\_lon}                 & 0.0340 \\
\textit{dinamana\_minutes}                & 0.0207 \\
\textit{vedic\_chandra\_speed\_deg/day}   & 0.0203 \\
\textit{spline\_bearing\_deg}             & 0.0187 \\
\textit{elevation}                        & 0.0179 \\
\textit{ratrimana\_minutes}               & 0.0174 \\
\textit{mars\_ecliptic\_lat\_deg}         & 0.0168 \\
\textit{sun\_dec\_deg}                    & 0.0162 \\
\bottomrule
\end{tabular}
\end{table}

\subsubsection{Caution on Spline Features}

The features \textit{spline\_end\_lon}, \textit{spline\_bearing\_deg},
\textit{spline\_curvature}, \textit{spline\_displacement}, and
\textit{spline\_length\_deg} encode the smoothed forecast track geometry.
Using these to predict current or next-step wind introduces potential
\emph{implicit leakage}: knowing the future trajectory implies knowing the
future position of the storm. These features rank highly in importance but
are excluded from the scientific (astronomical) block in the predictive model
to avoid contaminating the comparison.

\subsection{SHAP / LIME Explainability}
\label{sec:shap}

The SHAP analysis uses an XGBoost model trained on features with $> 80\%$
completeness (186 features) with a random 80/20 split (train $R^2 = 0.784$,
test $R^2 = 0.707$).
\textbf{Note:} This $R^2$ is inflated relative to the grouped-CV estimate
($\approx -0.10$) due to within-storm autocorrelation in the random split.

Fourteen features appear in the top-30 of \emph{both} SHAP and LIME importance
rankings (consensus set):

\begin{itemize}
  \item Geospatial: \textit{latitude}, \textit{longitude}, \textit{nearest\_lat}
  \item Track geometry (leakage risk): \textit{spline\_bearing\_deg}, \textit{spline\_end\_lon}
  \item Day-length: \textit{dinamana\_minutes}, \textit{ratrimana\_minutes}
  \item Solar: \textit{sun\_dec\_deg}
  \item Planetary: \textit{mercury\_dec\_deg}, \textit{neptune\_ecliptic\_lon\_deg}
  \item Vedic planetary speeds/degrees: \textit{vedic\_arun\_speed\_deg/day}
        (Uranus), \textit{vedic\_budha\_full\_degree} (Mercury),
        \textit{vedic\_budha\_speed\_deg/day}, \textit{vedic\_guru\_full\_degree} (Jupiter)
\end{itemize}

After removing leakage-risk spline features and day-length features
(which encode latitude-driven seasonality), the robust non-redundant signals
are \emph{latitude}, \emph{longitude}, and several planetary position quantities.
Planetary quantities with high individual SHAP importance
(e.g.\ \textit{vedic\_chandra\_speed\_deg/day}, \textit{mars\_ecliptic\_lat\_deg})
likely reflect a confound: planetary positions are correlated with calendar time
and thus with the tropical cyclone season.

\subsection{Hypothesis Testing Results}
\label{sec:hypotheses}

\subsubsection{H1: Vedic Predictive Lift}

We trained a Random Forest model to predict \textit{wind\_change} using:
(a) meteorological features only (baseline), and
(b) meteorological + Vedic features.

\begin{center}
\begin{tabular}{lcc}
\toprule
\textbf{Model} & \textbf{RMSE (kt)} & \textbf{Change} \\
\midrule
Meteorological baseline & 4.04 & --- \\
Meteorological + Vedic  & 4.55 & $+12.6\%$ (worse) \\
\bottomrule
\end{tabular}
\end{center}

\textbf{Conclusion (H1):} Vedic features degrade prediction. The null hypothesis
that Vedic features have no predictive lift is \emph{not} rejected; Vedic features
actively harm out-of-sample performance, consistent with overfitting noise.

\subsubsection{H2: Lunar Phase and Rapid Intensification}

We define Rapid Intensification (RI) as a wind increase of $\geq 30$~kt within
a 24-hour rolling window (note: WMO standard is $\geq 35$~kt; we document this
discrepancy). With this threshold, 2\,020 RI events were identified across
47\,533 observations.

A chi-square test of the contingency table (lunar phase $\times$ RI/non-RI):
\begin{equation}
  \chi^2 = 19.66, \quad \text{df} = 7, \quad p = 0.0064
\end{equation}

With Bonferroni correction across 7 main hypothesis groups
($\alpha_{\text{adj}} = 0.05/7 = 0.0071$), this result \emph{barely} survives
correction ($p = 0.0064 < 0.0071$).

The Waning Crescent phase shows the highest RI count (301 events vs.\ a
cross-phase mean of $\approx 253$). The effect size (Cram\'er's $V \approx 0.02$)
is small. \textbf{Conclusion (H2):} Marginal statistical evidence for a
lunar-phase -- RI association; physical interpretation requires caution.

\subsubsection{H3: Planetary Speeds and Trajectory Dynamics}

Spearman correlation was computed between 60 pairs of (planetary speed, storm
trajectory feature). The top associations all involve outer-planet speed proxies
(Jupiter, Saturn, Neptune) correlated with day-of-year encoded trajectory variables.
Retrograde vs.\ direct ANOVA showed no significant wind-change differences for
any tested body after Bonferroni correction.

\subsubsection{Samvatsara Analysis}

The 60-year Vedic Samvatsara cycle maps, in the current dataset, to approximately
one calendar year per Samvatsara (since $\approx 10$ years of data are available,
roughly 10 Samvatsara names appear). A one-way ANOVA of wind speed by Samvatsara
reaches $F$-stat $\gg 10^{100}$, $p < 10^{-130}$.

\textbf{Critical caveat:} This test is mathematically equivalent to asking
``does storm intensity differ across calendar years?'' The answer is trivially
yes due to interannual climate variability (ENSO, PDO, Atlantic Multidecadal
Oscillation). The result carries no evidence for a Hindu calendar-cycle effect.

\subsubsection{Mars / Rahu / Ketu Analysis}

\begin{center}
\begin{tabular}{lcc}
\toprule
\textbf{Analysis} & \textbf{Statistic} & \textbf{$p$-value} \\
\midrule
Mars+Rahu+Ketu model vs.\ baseline (RMSE) & $+1.1\%$ (worse) & --- \\
Mars altitude band $\times$ RI ($\chi^2$)  & --- & 0.532 \\
Retrograde vs.\ direct wind-change (Mars)  & ANOVA $F$ & 0.916 \\
Rahu retrograde vs.\ direct wind-change    & ANOVA $F$ & 0.040 \\
Rahu degree vs.\ wind-change (Spearman)    & $\rho$     & 0.006 \\
Ketu retrograde vs.\ direct wind-change    & ANOVA $F$ & 0.040$^*$ \\
\bottomrule
\multicolumn{3}{l}{\small $^*$ Identical to Rahu by astronomical constraint (see Section~\ref{sec:corrections}).}
\end{tabular}
\end{center}

Rahu degree Spearman $p = 0.006$ and retrograde ANOVA $p = 0.040$ fail Bonferroni
correction ($\alpha_{\text{adj}} = 0.007$). Mars, Rahu, and Ketu features do not
survive family-wise error rate control.

\subsection{Block-Wise Predictive Modelling}
\label{sec:models}

Table~\ref{tab:models} summarises the four feature blocks in temporal
forward-chaining cross-validation (train / test split: 2016--2023~Oct / 2023--2026).

\begin{table}[h]
\centering
\caption{Block-wise Random Forest regression results (predicting next-step wind).
  Lower RMSE is better.}
\label{tab:models}
\begin{tabular}{lcc}
\toprule
\textbf{Feature Block} & \textbf{Test RMSE (kt)} & \textbf{Test MAE (kt)} \\
\midrule
Meteorological only (14 features)    & \textbf{4.164} & \textbf{2.540} \\
Stacked residual                     & 4.193          & 2.583 \\
Meteo + Panchang (18 added)          & 4.372          & 2.685 \\
Full (all blocks, 190 added)         & 4.515          & 2.921 \\
\bottomrule
\end{tabular}
\end{table}

\textbf{Key finding:} The meteorological-only model achieves the best RMSE.
Each additional block of astronomical or Vedic features increases RMSE,
with the full model ($+190$ features) performing worst. This is consistent
with the hypothesis testing result (H1) and confirms that the 193 astronomical
features add noise.

\textbf{RI classification note:} The adaptive RI threshold selector for the
classification task settled on $\geq 10$~kt (next-step delta), which is
fundamentally different from the 30~kt/24h rolling definition used in H2.
Cross-script RI comparisons should not be made.

% ═══════════════════════════════════════════════════════════════════════════════
\section{Discussion}
% ═══════════════════════════════════════════════════════════════════════════════

\subsection{The Seasonal Confound}
\label{sec:confound}

The strongest Vedic categorical associations (Drik Ritu, Sunsign, Drik Ayana)
all encode time of year: tropical cyclone seasons peak in summer--autumn months,
so any calendar variable correlated with the season will show a wind association.
This is not a Vedic astronomy effect; it is the tropical cyclone season signal
re-expressed in calendar coordinates.

Similarly, Vedic planetary positions (full degree, speed) are functions of calendar
time for the slow outer planets (Saturn, Jupiter, Rahu). Their apparent correlation
with storm intensity is a proxy for year-level interannual variability.

\subsection{Multicollinearity Between Vedic and Astronomical Features}

Vedic planetary features (e.g.\ \textit{vedic\_budha\_full\_degree}) are computed
from the same Skyfield ephemerides as the Skyfield features
(e.g.\ \textit{mercury\_ecliptic\_lon\_deg}) but with different coordinate systems
(sidereal vs.\ tropical ecliptic, offset by the Lahiri ayanamsha $\approx 23.9\textdegree$).
Any observed Vedic feature correlation is therefore geometrically equivalent to
the corresponding Skyfield feature correlation. There are 239 feature pairs with
$|r| > 0.95$; adding Vedic features on top of Skyfield features is nearly redundant
and degrades models through noise amplification.

\subsection{Variance Inflation Factors}

VIF analysis on a representative 9-feature subset confirms low multicollinearity
(all VIF $\leq 2.0$, well below the conventional danger thresholds of 5 and 10),
suggesting the chosen meteorological predictors are not strongly collinear with
each other. The problem is that the 193 astronomical features are collectively
redundant \emph{with each other}, not with the meteorological features.

\subsection{Interpretation of SHAP / LIME Consensus Features}

Of the 14 consensus features, 7 are straightforward seasonal or geographic signals:
\textit{latitude}/\textit{longitude} (climatological intensity zone),
\textit{dinamana\_minutes}/\textit{ratrimana\_minutes} (day length = latitude + season),
\textit{sun\_dec\_deg} (solar declination = season).
Three are planetary quantities that, as argued above, likely capture seasonality.
Only after removing these geographic and seasonal proxies do the remaining features
(\textit{vedic\_budha\_speed\_deg/day}, \textit{vedic\_guru\_full\_degree},
\textit{vedic\_arun\_speed\_deg/day}) represent potential astrological signals;
and even these are weak (importance $< 0.02$) in a model that generalises poorly
across storms.

\subsection{Limitations}

\begin{enumerate}
  \item \textbf{No meteorological driver features}: The dataset does not include
        SST, vertical wind shear, moisture, or upper-level flow, which are the
        primary determinants of tropical cyclone intensity. Absence of these features
        explains the negative grouped CV $R^2$ for all models.
  \item \textbf{Coarse temporal resolution}: Panchang quantities and Skyfield
        ephemerides are computed at each track advisory time ($\approx 6$-hourly),
        but cyclone intensity change operates on sub-6-hour scales during rapid
        intensification.
  \item \textbf{Multiple basins}: The dataset covers all global basins simultaneously.
        Basin-specific effects (Atlantic AMO, Pacific ENSO sensitivity) are not
        separated, adding variance that confounds between-year comparisons.
  \item \textbf{Spline leakage}: Spline track features may encode future trajectory
        information; their high importance scores should not be interpreted as evidence
        of out-of-sample predictive power.
  \item \textbf{Single ML family}: Feature importance and SHAP are computed from
        tree-ensemble models only; linear and Gaussian process models might yield
        different feature rankings.
\end{enumerate}

% ═══════════════════════════════════════════════════════════════════════════════
\section{Conclusion}
% ═══════════════════════════════════════════════════════════════════════════════

We conducted a rigorous, multi-method exploratory investigation into whether
Vedic Panchang and astronomical ephemeris quantities carry statistically meaningful
associations with tropical cyclone intensity and dynamics.

Our principal findings are:
\begin{enumerate}
  \item Vedic features consistently \emph{degrade} out-of-sample predictive
        performance relative to a meteorological-only baseline ($+12.6\%$ RMSE in H1,
        $+8.5\%$ RMSE in the block-wise model). There is no evidence of predictive lift.
  \item The strongest Vedic category associations (Drik Ritu, Sunsign) are seasonal
        proxies; they encode the tropical cyclone season signal, not astrological effects.
  \item A marginal lunar-phase -- RI association ($p = 0.0064$) survives narrow
        Bonferroni correction but has Cram\'er's $V \approx 0.02$, implying negligible
        practical importance.
  \item All three tree-ensemble models achieve negative grouped-CV $R^2$ on the full
        feature set, confirming the astronomical feature block adds noise.
  \item The dominant predictive features remain geospatial coordinates
        (\textit{latitude}, \textit{longitude}) and track-geometry quantities
        (\textit{spline\_bearing\_deg}), consistent with established tropical
        cyclone climatology.
\end{enumerate}

This null result has practical value: it defines the upper bound of any Vedic /
Panchang contribution to cyclone intensity that can be detected in a decade-long
global dataset at 6-hourly resolution. Future work should (a) include physical
meteorological drivers (SST, wind shear), (b) focus basin-specific analyses,
(c) apply gravitational tidal forcing models as a physically motivated replacement
for empirical Vedic categories, and (d) explore whether lunar-phase tidal signals
appear more clearly in SST anomaly or mixed-layer depth data that preconditions
the ocean for intensification.

% ═══════════════════════════════════════════════════════════════════════════════
\section*{Acknowledgements}
% ═══════════════════════════════════════════════════════════════════════════════

The Skyfield ephemeris computations used the Skyfield Python library and the
DE421 JPL planetary ephemeris. Vedic Panchang calculations used a Drig-Ganita
(observational) Panchang computation engine. We thank the open-source communities
behind NumPy, pandas, scikit-learn, XGBoost, LightGBM, SHAP, LIME, matplotlib,
and seaborn.

% ═══════════════════════════════════════════════════════════════════════════════
% Bibliography
% ═══════════════════════════════════════════════════════════════════════════════
\bibliographystyle{plainnat}
\begin{thebibliography}{9}

\bibitem[DeMaria et~al.(2014)]{demaria2014advances}
DeMaria, M., Sampson, C.~R., Knaff, J.~A., and Musgrave, K.~D. (2014).
Is tropical cyclone intensity guidance improving?
\textit{Bulletin of the American Meteorological Society}, 95(3), 387--398.

\bibitem[Zhang \& Tao (2018)]{zhang2018tropical}
Zhang, J.~A. and Tao, D. (2018).
Impacts of vertical wind shear on the predictability of tropical cyclones:
a deep learning study.
\textit{Geophysical Research Letters}, 45(9), 4289--4297.

\bibitem[Rhodes (2019)]{skyfield2019}
Rhodes, B.~A. (2019).
Skyfield: High precision research-grade positions for planets and Earth satellites
generator.
\textit{Astrophysics Source Code Library}, record ascl:1907.024.

\bibitem[Kaplan \& DeMaria (2003)]{kaplan2003}
Kaplan, J. and DeMaria, M. (2003).
Large-scale characteristics of rapidly intensifying tropical cyclones in the
North Atlantic basin.
\textit{Weather and Forecasting}, 18(6), 1093--1108.

\bibitem[Benjamini \& Hochberg (1995)]{bh1995}
Benjamini, Y. and Hochberg, Y. (1995).
Controlling the false discovery rate: A practical and powerful approach to
multiple testing.
\textit{Journal of the Royal Statistical Society Series B}, 57(1), 289--300.

\bibitem[Lundberg \& Lee (2017)]{lundberg2017}
Lundberg, S.~M. and Lee, S.-I. (2017).
A unified approach to interpreting model predictions.
\textit{Advances in Neural Information Processing Systems}, 30, 4765--4774.

\bibitem[Ribeiro et~al.(2016)]{ribeiro2016}
Ribeiro, M.~T., Singh, S., and Guestrin, C. (2016).
``Why should I trust you?'': Explaining the predictions of any classifier.
\textit{Proceedings of KDD 2016}, 1135--1144.

\end{thebibliography}

% ═══════════════════════════════════════════════════════════════════════════════
% Appendices
% ═══════════════════════════════════════════════════════════════════════════════
\appendix
\section{Feature Block Definitions}
\label{app:blocks}

\begin{description}
  \item[Meteorological block (14 features)] \textit{wind}, \textit{pressure},
    \textit{latitude}, \textit{longitude}, \textit{wind\_change},
    \textit{pressure\_change}, \textit{displacement}, \textit{movement\_speed},
    \textit{hours\_elapsed}, \textit{spline\_bearing\_deg},
    \textit{spline\_curvature}, \textit{spline\_displacement},
    \textit{spline\_length\_deg}, \textit{elevation}.

  \item[Panchang block (18 features)] Tithi (numeric), Nakshatra (numeric),
    Yoga (numeric), Karana (numeric), Paksha (binary), Moonsign (numeric),
    Sunsign (numeric), Drik Ritu (numeric), Drik Ayana (numeric),
    Samvatsara (numeric), \textit{dinamana\_minutes}, \textit{ratrimana\_minutes},
    day-of-week, day-of-year, month, hour\_utc, \textit{lahiri\_ayanamsha}.

  \item[Scientific block (190 features)] All Skyfield per-body position features
    (RA/Dec, ecliptic lon/lat, altitude/azimuth, sin/cos transforms) plus
    Vedic planetary numeric quantities (full degree, padam, speed, RA, Dec).
    \emph{Excludes} \textit{spline\_} prefix features to avoid leakage.
\end{description}

\section{Circular-Linear Correlation Results (Top 10)}
\label{app:circular}

\begin{center}
\begin{tabular}{lccc}
\toprule
\textbf{Vedic degree feature} & $r_{\sin}$ & $r_{\cos}$ & $r_{\text{circ}}$ \\
\midrule
\multicolumn{4}{c}{\small (see \texttt{EDA/output/comprehensive/vedic\_degree\_correlations.csv} for full table)} \\
\bottomrule
\end{tabular}
\end{center}

All circular-linear correlations with wind speed are small ($|r_{\text{circ}}| < 0.15$
for all tested Vedic degree features), confirming that no angular Vedic planetary
feature has a meaningful linear relationship with cyclone intensity.

\section{Pipeline Audit Summary}
\label{app:audit}

\begin{center}
\begin{tabular}{llll}
\toprule
\textbf{ID} & \textbf{Script} & \textbf{Issue} & \textbf{Severity} \\
\midrule
A2 & 06 & Rahu = Ketu retrograde (astronomical constraint) & Critical \\
B1 & 02 & No multiple-comparison correction (fixed: BH FDR) & Medium \\
C1 & 03 & Negative CV $R^2$; feature ranks unreliable & Critical \\
C3 & 03 & Same-storm CV leakage (fixed: GroupKFold) & Medium \\
D1 & 04 & $R^2$ discrepancy vs.\ Stage~3 (documented) & Critical \\
E2 & 05 & Circular degrees treated as linear (fixed: sin/cos) & Critical \\
E3 & 05 & Retrograde no $p$-values (fixed: Mann--Whitney U) & Medium \\
F2 & 06 & Vedic model worse than baseline (flagged in report) & High \\
F5 & 06 & Samvatsara ANOVA = interannual confound (documented) & High \\
F6 & 06 & No Bonferroni correction (fixed: $\alpha_{\text{adj}} = 0.007$) & Medium \\
G3 & 07 & Spline feature leakage (removed from sci.\ block) & Critical \\
G2 & 07 & Best model = meteo\_only, not full (flagged in report) & High \\
\bottomrule
\end{tabular}
\end{center}

\end{document}
